\section{La materia vivente}

% ---------------------------------------------------------------------------
\subsection{Proprietà e componenti}

\rubrica{Proprietà della materia vivente}
\begin{itemize}
  \item metabolismo;
  \item accrescimento;
  \item moltiplicazione;
  \item irritabilità.
\end{itemize}

\textbf{Cellula}: la più piccola parte di protoplasma capace di un'esistenza indipendente; è l'unità morfologica, funzionale e fisiologica di tutti gli organismi viventi.

\rubrica{Componenti}
Il 75--85\% del protoplasma è costituito da acqua.
\begin{itemize}
  \item \textbf{inorganici}:
    \begin{itemize}
      \item elementi primari (97--98\%): C, H, O, N, S, P;
      \item cationi: $K^+$, $Na^+$, $Ca^{2+}$, $Mg^{2+}$;
      \item anioni: $Cl^-$, $HCO_3^-$, $PO_4^{3-}$, $SO_4^{3-}$;
      \item oligoelementi: Mn, Fe, Co, Cu, Mo, I, Br, V, Al, Se;
    \end{itemize}
  \item \textbf{organici}: carboidrati, lipidi, proteine, acidi nucleici.
\end{itemize}

% ---------------------------------------------------------------------------
\subsection{Glicidi o carboidrati}

Costituiscono più della metà dell'apporto energetico totale; sono formati da carbonio, idrogeno e ossigeno. Si dividono in monosaccaridi, disaccaridi/oligosaccaridi e polisaccaridi.

\rubrica{Monosaccaridi}
\begin{itemize}
  \item \textbf{triosi}: gliceraldeide;
  \item \textbf{pentosi}: ribosio;
  \item \textbf{esosi}: glucosio.
\end{itemize}
Note:
\begin{itemize}
  \item apporto giornaliero di carboidrati: 250--800 g nella dieta dei paesi occidentali;
  \item il glucosio è il nutriente fondamentale;
  \item esosi e pentosi: ruolo fondamentale nei processi nutritivi cellulari;
  \item glucosio e fruttosio: presenti in forma libera in alcuni frutti e vegetali.
\end{itemize}

\rubrica{Disaccaridi e oligosaccaridi}
Associazione di 2 o più molecole di monosaccaridi:
\begin{itemize}
  \item \textbf{saccarosio}: glucosio $+$ fruttosio;
  \item \textbf{maltosio}: glucosio $+$ glucosio;
  \item \textbf{lattosio}: glucosio $+$ galattosio.
\end{itemize}

\rubrica{Polisaccaridi}
Concatenazione di numerose unità monosaccaridiche:
\begin{itemize}
  \item \textbf{omopolisaccaridi}: unità monosaccaridiche identiche;
  \item \textbf{eteropolisaccaridi}: unità monosaccaridiche diverse.
\end{itemize}

\textbf{Amido}: principale polisaccaride di deposito nei vegetali. \textbf{Glicogeno}: principale polisaccaride di deposito negli animali. Entrambi costituiti da amilosio e amilopectina:
\begin{itemize}
  \item \textit{amilosio}: lineare, molecole di glucosio unite da legame monoglicosidico $\alpha(1\rightarrow4)$;
  \item \textit{amilopectina}: ramificata, legami glicosidici $\alpha(1\rightarrow6)$ nei punti di ramificazione.
\end{itemize}

\textbf{Cellulosa}: polimero lineare formato da migliaia di monomeri di glucosio uniti da legame $\beta(1\rightarrow4)$.
\begin{itemize}
  \item le catene si dispongono parallelamente mediante legami idrogeno, formando fibrille;
  \item funzione strutturale nel regno vegetale: costituisce la parete primaria e secondaria delle cellule vegetali;
  \item fermentata dalla flora batterica intestinale.
\end{itemize}

% ---------------------------------------------------------------------------
\subsection{Lipidi}

Due grandi classi:
\begin{itemize}
  \item \textbf{triacilgliceroli} (trigliceridi) --- 98\%: funzione di riserva energetica; lipidi di deposito, negli adipociti del tessuto adiposo; grassi animali;
  \item \textbf{fosfolipidi, glicolipidi, colesterolo, vitamine liposolubili} --- 2\%: funzione strutturale (membrane cellulari); dal colesterolo derivano la sintesi dei sali biliari e della vitamina D.
\end{itemize}

\rubrica{Acidi grassi}
Catene lineari con numero \textbf{pari} di atomi di carbonio; presentano un gruppo carbossilico (--COOH).

\textit{Saturi}:
\begin{itemize}
  \item catena corta $\rightarrow$ solubili in acqua e volatili: formico ($HCOOH$), acetico ($CH_3COOH$), propionico ($CH_3CH_2COOH$), butirrico ($CH_3CH_2CH_2COOH$);
  \item catena media: da 6 a 12 C (caprilico, laurico);
  \item catena lunga $\rightarrow$ insolubili in acqua: da 14 a 18 C (palmitico, stearico);
  \item catena molto lunga: da 20 a 24 C (arachidico).
\end{itemize}

\textit{Insaturi}: presenza di uno o più doppi legami; la solubilità in acqua aumenta con l'aumentare dei doppi legami:
\begin{itemize}
  \item monoinsaturo: acido oleico;
  \item polinsaturo: acido linoleico.
\end{itemize}

\rubrica{Lipidi semplici}
Derivano dall'esterificazione degli acidi grassi con alcoli:
\[
\text{glicerolo} + \text{acido grasso} \rightarrow \text{monogliceride} + \text{acqua}.
\]
\textbf{Sintesi dei trigliceridi}: glicerolo $+$ 3 acidi grassi $\rightarrow$ trigliceride $+$ 3 $H_2O$.

\textbf{Triacilgliceroli}: presenti nel tessuto adiposo; funzioni: riserva energetica, isolamento termico, sostegno di organi.

\rubrica{Lipidi complessi}
\textbf{Fosfolipidi}: molecole anfipatiche:
\begin{itemize}
  \item testa \textbf{polare}: gruppo fosfato (con glicerolo);
  \item code \textbf{apolari}: catene di acidi grassi.
\end{itemize}
Esempi: fosfatidilcolina (PC), fosfatidilinositolo (PI), fosfatidilserina (PS), fosfatidiletanolammina (PE), sfingomielina (SM).

\rubrica{Steroidi}
Lipidi policiclici la cui struttura di base è data da quattro anelli fusi tra loro. Esempi: colesterolo, vitamina D.

% ---------------------------------------------------------------------------
\subsection{Amminoacidi}

Struttura generale: un carbonio centrale ($\alpha$) legato a quattro gruppi:
\begin{itemize}
  \item gruppo carbossilico (--COOH);
  \item gruppo amminico (--$NH_2$);
  \item gruppo laterale R, proprio di ciascun amminoacido;
  \item un atomo di idrogeno.
\end{itemize}
La natura delle catene laterali conferisce proprietà diverse alle proteine in cui sono presenti.

\rubrica{Classificazione in base alla catena laterale}
\begin{itemize}
  \item \textbf{non polari}: alanina, valina, leucina, metionina, isoleucina, prolina, fenilalanina, triptofano;
  \item \textbf{polari ma non carichi}: cisteina, glicina, serina, treonina, asparagina, glutammina, tirosina;
  \item \textbf{carichi negativamente}: glutammato, aspartato;
  \item \textbf{carichi positivamente}: arginina, lisina, istidina.
\end{itemize}

\rubrica{Classificazione nutrizionale}
\begin{itemize}
  \item \textbf{essenziali} (l'organismo umano non è in grado di sintetizzarli, essenziali in ogni circostanza): fenilalanina, isoleucina, leucina, lisina, metionina, triptofano, treonina, valina, istidina;
  \item \textbf{condizionatamente essenziali} (richiedono un apporto esogeno, la sintesi endogena non è sufficiente): arginina, cisteina, glutammina, glicina, prolina, tirosina;
  \item \textbf{non essenziali} (sintetizzati dall'organismo in quantità sufficienti): alanina, aspartato, asparagina, glutammato, serina.
\end{itemize}

% ---------------------------------------------------------------------------
\subsection{Peptidi e legame peptidico}

\textbf{Legame peptidico}: si forma dalla condensazione del gruppo carbossilico (--COOH) del primo amminoacido con il gruppo amminico (--$NH_2$) del secondo, con formazione del legame --CO--NH-- e rilascio di una molecola d'acqua (reazione di sintesi disidratativa).
\begin{itemize}
  \item estremità \textbf{N-terminale}: gruppo amminico libero;
  \item estremità \textbf{C-terminale}: gruppo carbossilico libero.
\end{itemize}

% ---------------------------------------------------------------------------
\subsection{Proteine}

Quattro livelli di organizzazione strutturale.

\rubrica{Struttura primaria}
Data dalla sequenza amminoacidica; specifica di ogni proteina.

\rubrica{Struttura secondaria}
Ripiegamenti regolari della catena, mantenuti da legami idrogeno:
\begin{itemize}
  \item \textbf{$\alpha$-elica}: avvolgimento a spirale; i legami idrogeno mantengono l'avvolgimento; i gruppi laterali si proiettano all'esterno dai lati;
  \item \textbf{foglietto $\beta$}: la catena polipeptidica si ripiega su se stessa; metà dei gruppi laterali proiettata al di sopra del foglietto, l'altra metà al di sotto; i filamenti possono disporsi in direzioni opposte (antiparallelo) o nella stessa direzione (parallelo).
\end{itemize}

\rubrica{Struttura terziaria}
Organizzazione tridimensionale complessiva; dipende dalle interazioni tra i gruppi laterali:
\begin{itemize}
  \item legame a idrogeno;
  \item legame ionico;
  \item ponte disolfuro (S--S);
  \item interazione idrofobica.
\end{itemize}
Con la \textbf{denaturazione} la proteina perde la propria organizzazione ($\alpha$-eliche e filamenti $\beta$); durante la \textbf{rinaturazione} i ponti disolfuro aiutano a recuperare la forma nativa.

\rubrica{Struttura quaternaria}
Formata dall'unione di più catene polipeptidiche. Esempi: emoglobina, collagene.

\rubrica{Funzioni delle proteine}
\begin{itemize}
  \item \textbf{ormoni}: veicolano segnali regolatori tra le cellule;
  \item \textbf{anticorpi}: difesa;
  \item \textbf{accumulo}: riserva di molecole;
  \item \textbf{strutturali}: sostegno;
  \item \textbf{enzimatiche}: catalizzatori;
  \item \textbf{trasporto}: proteine canale;
  \item \textbf{motrici}: movimento cellulare;
  \item \textbf{regolatori}: regolano l'attività di altre molecole;
  \item \textbf{recettori}: ricevono segnali dall'ambiente extracellulare.
\end{itemize}

% ---------------------------------------------------------------------------
\subsection{Acidi nucleici}

Due tipi:
\begin{itemize}
  \item \textbf{DNA}: acido desossiribonucleico;
  \item \textbf{RNA}: acido ribonucleico.
\end{itemize}

\rubrica{Nucleosidi e nucleotidi}
\begin{itemize}
  \item \textbf{nucleoside}: zucchero (pentoso) $+$ base azotata, uniti da legame glicosidico;
  \item il pentoso è \textit{ribosio} (--OH in $2'$) o \textit{desossiribosio} (--H in $2'$);
  \item \textbf{nucleotide}: nucleoside $+$ gruppo fosfato $\rightarrow$ nucleoside monofosfato (NMF), difosfato (NDF), trifosfato (NTF).
\end{itemize}

\rubrica{Basi azotate}
\begin{itemize}
  \item \textbf{puriniche}: adenina, guanina;
  \item \textbf{pirimidiniche}: citosina, uracile, timina.
\end{itemize}
Appaiamento tra basi complementari mediante legami idrogeno: A--T e G--C nel DNA; A--U nell'RNA.

% ---------------------------------------------------------------------------
\subsection{Il DNA}

\rubrica{Struttura}
\begin{itemize}
  \item doppia elica; scheletro fosfato--desossiribosio, con legame fosfodiesterico;
  \item i due filamenti sono antiparalleli ($5'\rightarrow3'$);
  \item larghezza 2,2 nm; passo dell'elica $=$ 10 paia di basi $=$ 3,4 nm (20\,\AA{} di larghezza, 34\,\AA{} di passo, 3,4\,\AA{} per base).
\end{itemize}
Struttura proposta da Watson e Crick (\textit{Nature} 171, 737--738, 25 aprile 1953); l'analisi del DNA mediante diffrazione dei raggi X fu ottenuta da Rosalind Franklin.

\rubrica{Tre funzioni del materiale genetico}
\begin{enumerate}
  \item il DNA è depositario dell'informazione che codifica i caratteri ereditari;
  \item il DNA contiene l'informazione per dirigere la propria duplicazione;
  \item il DNA contiene l'informazione per dirigere la costruzione di proteine specifiche.
\end{enumerate}

\rubrica{Organizzazione della cromatina}
$DNA \rightarrow$ nucleosoma $\rightarrow$ fibra da 10 nm $\rightarrow$ fibra da 30 nm (solenoide) $\rightarrow$ domini ad ansa $\rightarrow$ impalcatura proteica $\rightarrow$ cromosoma metafasico.
\begin{itemize}
  \item \textbf{nucleosoma}: DNA avvolto attorno a un core di 8 istoni (2 molecole ciascuno di H2A, H2B, H3, H4);
  \item l'istone \textbf{H1} è associato al DNA linker tra i nucleosomi.
\end{itemize}

\rubrica{Replicazione}
\begin{enumerate}
  \item i due filamenti si svolgono e si separano;
  \item ciascun filamento funge da stampo per l'aggiunta di basi complementari (A--T, G--C);
  \item si formano 2 nuove doppie eliche identiche a quella parentale.
\end{enumerate}
Ogni molecola risultante è metà ``vecchia'' e metà ``nuova'' $\rightarrow$ replicazione \textbf{semiconservativa}.

\rubrica{DNA superavvolto}
Il DNA può presentarsi in forma rilassata o superavvolta (fortemente condensata). Esempio di compattazione: il DNA rilasciato dalla testa del batteriofago T2 è lungo 68\,$\mu$m, circa 60 volte la testa del fago nella quale è contenuto.

% ---------------------------------------------------------------------------
\subsection{L'RNA}

Caratteristiche:
\begin{itemize}
  \item singola catena nucleotidica;
  \item zucchero: ribosio;
  \item base: uracile (al posto della timina).
\end{itemize}

Diversi tipi:
\begin{itemize}
  \item \textbf{mRNA} (messaggero);
  \item \textbf{tRNA} (transfer): struttura a trifoglio; presenta l'\textit{anticodone} e, all'estremità $3'$, il sito di attacco dell'amminoacido;
  \item \textbf{rRNA} (ribosomiale): componente delle subunità ribosomiali.
\end{itemize}
